% !TEX root = ../../main.tex
% C5.4 — Thiết kế giao diện lập trình ứng dụng (API) (RÚT GỌN 2026-09-29: bỏ bảng mã trạng thái HTTP và bảng nhóm API; bản cũ ở report/archive/)
% Khớp backend/src/person_search/api: tiền tố /api/v1 (api/__init__.py), 9 blueprint (api/v1/*.py) + /health (live, ready, storage);
% error_response {"error": {code, message, details, request_id}} (api/errors.py); X-Request-ID; 429 + Retry-After (api/rate_limit.py);
% version trong PATCH users/cameras/cases -> 409 version_conflict; Idempotency-Key khi tải video; tải video trả 202;
% ảnh kết quả tìm kiếm /search-results/{id}/crop|frame, ảnh kết quả đã lưu /cases/{id}/results/{id}/crop|frame.
\section{Thiết kế giao diện lập trình ứng dụng}
\label{sec:5.4}

Ứng dụng web và máy chủ ứng dụng giao tiếp hoàn toàn qua một giao diện lập trình ứng dụng (API) theo phong cách REST. Mọi API nghiệp vụ nằm dưới tiền tố \path{/api/v1} và gọi tên tài nguyên bằng danh từ số nhiều, như \path{/cases}. Các API chỉ dành cho Quản trị viên nằm dưới \path{/admin}. Dữ liệu trao đổi ở dạng JSON, riêng ảnh truy vấn và tệp video được gửi dạng \texttt{multipart/form-data}. Các thao tác chuyển trạng thái có ý nghĩa nghiệp vụ, như khóa tài khoản hay loại camera khỏi vận hành, được thiết kế thành hành động tường minh, ví dụ \texttt{POST} \path{/admin/cameras/{id}/retire}, để mỗi thao tác được kiểm tra điều kiện và ghi nhật ký riêng.

Mọi lỗi được trả về theo một định dạng thống nhất gồm \emph{mã lỗi} ổn định để giao diện xử lý theo chương trình, \emph{thông báo} tiếng Việt để hiển thị, \emph{chi tiết} như lỗi của từng trường, và \emph{mã yêu cầu} để đối chiếu với nhật ký máy chủ. Mã trạng thái HTTP được dùng theo ý nghĩa chuẩn: 401 khi chưa đăng nhập hoặc phiên hết hạn, 403 khi vai trò không được phép, 404 khi tài nguyên không tồn tại hoặc nằm ngoài phạm vi được phép, 409 khi xung đột với trạng thái hiện tại, 422 khi dữ liệu không hợp lệ và 429 khi vượt giới hạn số lần gọi.

Để chống ghi đè, yêu cầu cập nhật tài khoản, camera và vụ việc phải kèm số phiên bản mà giao diện đang có. Nếu dữ liệu đã bị thay đổi ở nơi khác, máy chủ trả mã 409 thay vì ghi đè. Yêu cầu tải video có thể kèm một khóa chống lặp, nên gửi lại không tạo thêm công việc. Tải video lên cũng là một tác vụ không đồng bộ: máy chủ chỉ tạo công việc và trả mã 202 ngay, còn giao diện theo dõi tiến độ qua API công việc thay vì giữ yêu cầu mở trong suốt quá trình phân tích.

Về phân quyền, không API nào nhận khu vực của Giám sát viên hay người phụ trách vụ việc làm tham số. Máy chủ tự xác định các giá trị này từ phiên đăng nhập, nên việc sửa yêu cầu ở trình duyệt không thể mở rộng quyền. Phản hồi tìm kiếm và phản hồi xem vụ việc chỉ chứa thông tin mô tả, còn ảnh được lấy qua các đường dẫn riêng. Ảnh của kết quả tìm kiếm và ảnh của kết quả đã lưu dùng hai bộ đường dẫn khác nhau, ứng với hai căn cứ quyền xem ảnh ở Mục~\ref{sec:4.6}.

Ngoài các API nghiệp vụ, máy chủ có ba đường dẫn kiểm tra tình trạng (\path{/health/live}, \path{/health/ready}, \path{/health/storage}) phục vụ việc triển khai và giám sát.
